\clearpage\begingroup
\providecommand{\ObsSec}[2]{#1\mathord{\overset{\mathrm{s}}{,}}#2}
\providecommand{\ObsArc}[2]{#1\mathord{\overset{\prime\prime}{,}}#2}
\providecommand{\ObsJupiter}{\mathord{\tikz[baseline=-.1ex,x=1em,y=1em,line width=.045em]{
  \draw (.02,.48) .. controls (.04,.70) and (.33,.69) .. (.31,.46) .. controls (.29,.29) and (.14,.19) .. (.02,.13)--(.49,.13);
  \draw (.37,.64)--(.37,-.06);}}}
\providecommand{\ObsVesta}{\mathord{\tikz[baseline=-.1ex,x=1em,y=1em,line width=.045em]{
  \draw (.04,.02) rectangle (.40,.38);
  \draw (.22,.38)--(.22,.48) (.12,.63)--(.22,.48)--(.32,.63);}}}

\hypertarget{native-276-5}{}
\HistoricalHead{271}{R}{VERSCHIEDENE BEOBACHTUNGEN IN GÖTTINGEN.}
\def\ObsMoon{\mathord{\tikz[baseline=-.1ex,x=1em,y=1em,line width=.04em]{
  \draw (.42,.69) .. controls (-.06,.68) and (-.06,.05) .. (.42,.04);
  \draw (.42,.69) .. controls (.10,.59) and (.10,.14) .. (.42,.04);}}}
\def\ObsCeresR{\mathord{\tikz[baseline=-.1ex,x=1em,y=1em,line width=.045em]{
  \draw (.10,.52) .. controls (.22,.76) and (.47,.64) .. (.44,.43) .. controls (.42,.25) and (.23,.22) .. (.13,.33);
  \draw (.19,.34)--(.19,-.05);
  \draw (.05,.10)--(.33,.10);}}}
\noindent [5. Seite] 98, 1810 August 23.

\{Austritt Nr. 130 Tauri Flamst. am dunkeln $\ObsMoon$-Rande
\begin{center}\begin{tabular}{rrrll}
$23^{\mathrm h}$&$32^{\mathrm m}$&$\ObsSec{37}5$&\textsc{Gauss}&10 f. Herschel\\
&32&$37{,}8$&\textsc{Harding}&5 f. Geffken
\end{tabular}\end{center}
\noindent bei völlig heiterm Himmel. Sehr gute Beob.

Beim Eintritt am hellen $\ObsMoon$-Rande zitterte der $\ObsMoon$ zu sehr in den Dünsten des Horizonts, so dass der Stern kurz vor dem Eintritte unsichtbar ward.
\begin{center}\begin{tabular}{rr rrr l}
Wahrer Austritt in $\ast$ Zeit&$=$&$23^{\mathrm h}$&$37^{\mathrm m}$&$\ObsSec{45}{36}$&[G]\\
Mittl. Zeit&$=$&13&31&$10{,}46$&G\\
&&&&$10{,}8$&H\}\textsuperscript{[*)]}.
\end{tabular}\end{center}
\medskip\noindent [6. Seite] 109, 1810 September 18.
\begin{center}\small\setlength{\tabcolsep}{3pt}
\begin{tabular}{lrll}
\{Eintritt $\alpha$ Tauri am hellen $\ObsMoon$-Rande&$22^{\mathrm h}19^{\mathrm m}\ObsSec{25}7$&H[\textsc{arding}]&5 f. Reflect.\\
&$26{,}2$&G[\textsc{auss}]&10 f. Reflect.
\end{tabular}
\end{center}
\noindent\hspace*{\parindent}Austritt $\alpha$ Tauri am dunkeln $\ObsMoon$-Rande $23^{\mathrm h}14^{\mathrm m}\ObsSec{49}0$ G. und H.\}\textsuperscript{[*)]}.
\begin{center}\small\setlength{\arraycolsep}{3pt}
$\begin{array}{lrrrl}
\text{Sternzeit Eintritt}&22^{\mathrm h}&25^{\mathrm m}&\ObsSec{41}{07}&\text{\textsc{Harding}}\\
&&25&41{,}37&\text{\textsc{Gerling}}\\
&&25&41{,}57&\text{\textsc{Gauss}}\\
\hfill\text{Austritt}&23&21&4{,}51&\mathrm{H.\ G.\ G.}
\end{array}\quad
\left.\begin{array}{rrr}
10^{\mathrm h}&37^{\mathrm m}&\ObsSec43\\
&&4{,}6\\
&&4{,}8\\
11&32&18{,}7
\end{array}\right\}\begin{array}{l}
\text{\textsuperscript{[**)]}}\\
\mathrm{M.\ Z.}\\
\phantom{\mathrm{M.\ Z.}}\\
\phantom{\mathrm{M.\ Z.}}
\end{array}$
\end{center}
\begingroup\renewcommand{\thefootnote}{[*)}\footnotetext{\textsc{Harding}’s Handschrift.]}\endgroup
\begingroup\renewcommand{\thefootnote}{[**)}\footnotetext{\textsc{Gauss}’ Handschrift.]}\endgroup
\medskip\noindent [7. Seite] 126, 1811 Febr. 21.

Heute Abends wurde Ceres zweimal mit zwei Sternen der \textit{Hist. Cél.} verglichen, die dort pag. 65 so beobachtet sind:
\clearpage\SourceLeaf{277}\HistoricalHead{272}{L}{PRAKTISCHE UND SPHÄRISCHE ASTRONOMIE. NACHLASS UND BRIEFWECHSEL.}
\begin{center}\small\setlength{\tabcolsep}{4pt}
\begin{tabular}{lr|c|c|rrr|c|rrr|}
1794 April 25&7&$a$&—&$10^{\mathrm h}$&$28^{\mathrm m}$&$\ObsSec{45}3$&$29'16''$&$22^0$&$3'$&$25''$\\
&9&$b$&&&29&$50{,}5$&&22&5&5
\end{tabular}.
\end{center}
\noindent Meine Beobachtungen waren folgende:
\begin{center}\small\setlength{\tabcolsep}{5pt}
\begin{tabular}{rr|rr|rr|}
\multicolumn{2}{c}{$a$}&\multicolumn{2}{c}{$\ObsCeresR$}&\multicolumn{2}{c}{$b$}\\
$5^{\mathrm h}47^{\mathrm m}21^{\mathrm s}$&$50^{\mathrm m}7^{\mathrm s}$&$5^{\mathrm h}48^{\mathrm m}30^{\mathrm s}$&$56^{\mathrm m}56^{\mathrm s}$&$48^{\mathrm m}24^{\mathrm s}$&$51^{\mathrm m}\ObsSec{14}8$\\
$51\quad58$&$55\quad4$&$52\quad55$&$56\quad5$&$53\quad8$&$56\quad6$
\end{tabular}.
\end{center}
\medskip\noindent [8. Seite] 127, 1811 März 1\textsuperscript{[*)]}.

Der Austritt des $\alpha$ Tauri hinterm hellen Mondsrande
\begin{center}\begin{tabular}{l lrrr}
$8^{\mathrm h}17^{\mathrm m}15^{\mathrm s}$ Uhrzeit&Sternzeit&$8^{\mathrm h}$&$22^{\mathrm m}$&$\ObsSec{19}{89}$\\
&M. Z.&9&47&$16{,}4$.
\end{tabular}\end{center}
\medskip\noindent [9. Seite] 127, 1811 März 7.\textsuperscript{[*)]}.
\begin{center}\small\setlength{\tabcolsep}{4pt}\renewcommand{\arraystretch}{1.2}
\begin{tabular}{l|l|r|r|}
Eintritt o Leonis am&&Sternzeit&M. Z.\\
\quad dunkeln Rande&$10^{\mathrm h}36^{\mathrm m}41^{\mathrm s}$&$10^{\mathrm h}41^{\mathrm m}\ObsSec{45}{52}$&$11^{\mathrm h}42^{\mathrm m}\ObsSec{43}8$\\
Austritt&$11\quad47\quad7$ vielleicht $1^{\mathrm s}$ zu spät&$11\quad52\quad11{,}52$&$12\quad52\quad58{,}2$
\end{tabular}
\end{center}
\begingroup\renewcommand{\thefootnote}{[*)}\footnotetext{Vgl. auch Brief von \textsc{Gauss} an \textsc{Schumacher} vom 10. März 1811 und an \textsc{Olbers} vom 12. März 1811; der letztere wurde erst 1921 in einem Buche der Kiewer Sternwartenbibliothek aufgefunden und von \textsc{Grawe} im Jahresbericht der Neurussischen Naturforschenden Gesellschaft 1921 veröffentlicht; er wird auch in diesem Bande in der Abteilung Varia abgedruckt werden.]}\endgroup
\clearpage\endgroup
